\scanpage{50}
\mspage{\handwritten{- 59 -}}
% archive sheet number: 49
n\handwritten{°357}

p-adique $\handwritten{|}x\handwritten{|}_{p}=p^{-\unsure{v|\nu}_{p}(x)}$ (\uline{Top.Gén.} Chap.IX,\handwritten{§3}, n\handwritten{°2}) ; soit $\mathbb{Z}_{p}$ l'anneau des entiers p-adiques, adhérence de $\mathbb{Z}$ dans $\mathbb{Q}_{p}$ ; $\mathbb{Z}_{p}$ est un sous-groupe du groupe additif de $\mathbb{Q}_{p}$, et c'est aussi la boule ouverte $\handwritten{|}x\handwritten{|}_{p}\handwritten{<}3/2$ ; c'est donc un sous-groupe ouvert et fermé de $\mathbb{Q}_{p}$, de sorte que $\mathbb{Q}_{p}\handwritten{/}\mathbb{Z}_{p}$ est discret. De plus, $\mathbb{Z}_{p}$ est compact, par exemple parce qu'on peut l'identifier d'une manière évidente avec la limite projective des groupes $\mathbb{Z}\handwritten{/}p^{n}\mathbb{Z}$. Il est clair qu'il y a, sur le groupe additif $\mathbb{Q}_{p}$, une mesure de Haar $\handwritten{\alpha}$ et une seule telle que $\handwritten{\alpha}(\mathbb{Z}_{p})=1$ ; celle-ci est dite la \uline{mesure de Haar normalisée} sur $\mathbb{Q}_{p}$. Comme $\mathbb{Z}_{p}$ se compose, pour tout $n\handwritten{\geq}0$, de $p^{n}$ classes suivant le sous-groupe (ouvert et compact) $p^{-n}\mathbb{Z}_{p}$, on a $\handwritten{\alpha}(p^{n}\mathbb{Z}_{p})=p^{-n}$ ; on a de même $\handwritten{\alpha}(p^{-n}\mathbb{Z}_{p})=p^{n}$, puisque $p^{-n}\mathbb{Z}_{p}$, pour $n\handwritten{\geq}0$, se compose de $p^{n}$ classes suivant $\mathbb{Z}_{p}$. Soient $A_{1},\ldots,A_{m}$ des classes suivant $p^{n}\mathbb{Z}_{p}$ ; soit $f=\handwritten{\sum}_{i=1}^{m}a_{i}\handwritten{\unsure{\varphi|\phi}}_{A_{i}}$ ; on aura donc $\handwritten{\int}f\mathrm{d}\handwritten{\alpha}=p^{-n}\handwritten{\sum}_{i}a_{i}$. Soit maintenant $g\handwritten{\in}\handwritten{\unsure{\mathscr{K}|\mathcal{K}}}(\mathbb{Q}_{p})$ ; soit $K$ une réunion finie de classes suivant $\mathbb{Z}_{p}$, contenant le support de $g$. Comme $g$ est uniformément continue, il y a, pour tout $\handwritten{\varepsilon> }0$, un $n\handwritten{\geq}0$ tel que $\handwritten{|}g(x)-g(y)\handwritten{|\leq\varepsilon}$ pour $\handwritten{|}x-y\handwritten{|}_{p}\handwritten{\leq}p^{-n}$, c'est-à-dire pour $x\handwritten{\equiv}y\ \mathrm{mod.}\ p^{n}$ ; soient $A_{1},\ldots,A_{m}$ les classes suivant $p^{n}\mathbb{Z}_{p}$ contenues dans $K$ ; soient $a_{1},\ldots,a_{m}$ des représentants de ces classes ; soit $f=\handwritten{\sum}_{i=1}^{m}g(a_{i})\handwritten{\varphi}_{A_{i}}$ ; on a alors $\handwritten{|}f-g\handwritten{|\leq\varepsilon}$. En particulier, si $f$ ne prend qu'un nombre fini de valeurs distinctes $c_{\unsure{o|0}}=0$, $c_{1},\ldots,c_{r}$, et qu'on prenne $\handwritten{\varepsilon<}\inf_{i\handwritten{\ne}j}\handwritten{|}c_{i}-c_{j}\handwritten{|}$, on a $f=g$ ; plus particulièrement, si $f$ est la fonction caractéristique $\handwritten{\varphi}_{X}$ d'une ensemble $X$, et que $\handwritten{\varphi}_{X}\handwritten{\in\mathscr{K}}(\mathbb{Q}_{p})$, ce qui revient à dire que $X$ est ouvert et compact, on voit qu'il y a $n\handwritten{\geq}0$ tel que $X$ soit réunion des classes suivant $p^{n}\mathbb{Z}_{p}$ qui y sont contenues. Soit $\handwritten{\unsure{\Phi|\mathfrak{P}}}$ la famille formée de telles réunions, c'est-à-dire des parties ouvertes et compactes de $\mathbb{Q}_{p}$ ; c'est un clan au sens du Chap.IV,\handwritten{§4}, n\handwritten{°8}, et les fonctions $\handwritten{\Phi}$-étagées sont, d'après ce qu'on vient de voir, les fonctions
